The channel scale is defined using the adopted unit-cell hydraulic-diameter convention,
\begin{equation}
D_e=\frac{4A_0 L}{S},\qquad S=S_1+S_2 ,
\label{eq:De}
\end{equation}
where $A_0$ is the inlet flow area used in the saved cases, $L=l_c=l_d=25$~mm is the inlet-to-outlet cell length, and $S_1$ and $S_2$ are the two wetted, heated wall areas. This is an effective passage-scale definition: replacing the wetted perimeter in the usual $4A/P$ expression by $S/L$ gives Equation~\eqref{eq:De} \citep{zohuri2017,thulukkanam2013}. It should not be interpreted as a local hydraulic diameter at every cross-section of the embossed passage. Its numerical value is evaluated separately from each saved mesh.

The Reynolds number follows the standard definition \citep{cengel2014,thulukkanam2013},
\begin{equation}
\mathrm{Re}=\frac{uD_e}{\upsilon},\qquad \upsilon=\frac{\mu}{\rho}.
\label{eq:Re}
\end{equation}
The prescribed reference speed is therefore $u=4282\mu/(\rho D_e)$. Equal Reynolds number does not imply equal speed when $D_e$ varies: the saved cases give $u=0.308$--$0.679$~m/s. This is the reference mean speed used to scale the recycled velocity-component profiles, not an assumption of a uniform inlet profile or a substitution for the inlet-normal bulk velocity.

Fanning friction factor is adopted in this study, denoted $f_F$ (or simplified as $f$ in the plots). \citep{cengel2014,thulukkanam2013}
\begin{equation}
f_F=\frac{\Delta p\,D_e}{2\rho L u^2}.
\label{eq:f}
\end{equation}
Here $\Delta p$ is the difference between the area-averaged static pressures at the inlet and outlet. Thus the factor 2 in the denominator is intentional for the Fanning convention. All walls are stationary, no-slip and hydraulically smooth.

For constant $C_p$ and balanced inlet and outlet mass flow, the net fluid energy gain follows the steady-flow calorimetric balance \citep{thulukkanam2013,jiang2019a},
\begin{equation}
Q_f=\dot m C_p\left(T_{f,out}^{(E)}-T_{f,in}^{(E)}\right).
\label{eq:Q}
\end{equation}
Here $b\in\{in,out\}$ identifies the inlet or outlet boundary. For constant $C_p$, the energy-balance temperature is the signed temperature flow rate divided by the net mass flow \citep[Section~25.3.1, Equations~25--16 and~25--17]{ansys2025surface}:
\[
T_{f,b}^{(E)}=\frac{\sum_{i\in b}\dot m_i T_i}{\sum_{i\in b}\dot m_i},
\]
where $\dot m_i$ is the signed face mass flow rate relative to the boundary's outward normal; the net boundary mass flow is nonzero. Where local backflow occurs, conservation is checked using signed enthalpy-flux integrals rather than substituting absolute-flux temperature averages into the calorimetric balance \citep[Section~25.1]{ansys2025surface}. The wall heat rate used below is obtained directly from both heated walls, independently of this diagnostic.

For the reconstructed Nu results in Figure~\ref{figure4} and PEC comparisons in Figures~\ref{figure13}--\ref{figure15}, the case-specific reporting reference is the mass-weighted volume average of static temperature over the entire fluid domain $\Omega_f$, as defined in \citet[Section~25.4.1, Equation~25--39]{ansys2025surface}:
\[
T_{\mathrm{ref}}=T_{f,ave}
=\frac{\int_{\Omega_f}\rho T\,\mathrm{d}V}{\int_{\Omega_f}\rho\,\mathrm{d}V}
=\frac{\sum_{i\in\Omega_f}\rho_i V_i T_i}{\sum_{i\in\Omega_f}\rho_i V_i}.
\]
Here $V_i$, $\rho_i$ and $T_i$ are the volume, density and static temperature of fluid cell $i$. The integral is normalized by the total fluid mass, so $T_{f,ave}$ has units of temperature. It weights the mass contained in each cell, not the mass flow crossing a boundary. For the constant-density water used here, it is identical to the volume-weighted average $\sum_i V_i T_i/\sum_i V_i$. This fluid-domain average is distinct from the signed inlet and outlet temperatures $T_{f,b}^{(E)}$ used in the energy balance; it is not their arithmetic mean or difference.

The average heat-transfer coefficient is obtained by area-weighted averaging of Fluent's built-in \emph{Surface Heat Transfer Coef.} over the union of \texttt{wall\_1} and \texttt{wall\_2}, with the Reference Values temperature set to the case-specific $T_{f,ave}$. The equivalent value is reconstructed directly from the saved wall heat fluxes without modifying the original solution files. For the uniform wall temperature $T_w=333$~K, the Fluent field definition and surface-averaging procedure \citep{ansys2025htc} give
\begin{equation}
h_{ave}
=\frac{1}{S_1+S_2}\int_{S_1\cup S_2}
\frac{q''}{T_w-T_{\mathrm{ref}}}\,\mathrm{d}S
=\frac{Q_1+Q_2}{(S_1+S_2)(T_w-T_{f,ave})},
\qquad Q_j=\int_{S_j}q''\,\mathrm{d}S .
\label{eq:h_ave}
\end{equation}
Heat flux is positive from the heated walls into the water. Both wall heat rates and both wall areas therefore enter Equation~\eqref{eq:h_ave}; the individual wall coefficients are not added. This fluid-temperature-based form is consistent with the form of the mean coefficient in \citet{jiang2019a}, while the specific choice of a fluid-domain mass-average reference is defined explicitly here and is not attributed to that paper.

For $\theta_c=0\degree$, the fluid-domain mass-average static temperature is $T_{\mathrm{ref}}=T_{f,ave}=296.1537$~K, giving $h_{ave}=5346.87$~W\,m$^{-2}$\,K$^{-1}$. Across all 29 cases, $T_{f,ave}$ ranges from 295.4045 to 296.7044~K; each case uses its own value. The physical inlet temperature is 293~K and both walls are prescribed at 333~K. The reporting reference does not impose an additional thermal boundary condition or change the solved temperature field, wall heat rates or flow field.

The Nusselt and molecular Prandtl numbers use the standard nondimensional definitions \citep{thulukkanam2013},
\begin{equation}
\mathrm{Nu}=\frac{h_{ave}D_e}{\lambda},
\label{eq:Nu}
\end{equation}
\begin{equation}
\mathrm{Pr}=\frac{\mu C_p}{\lambda}=\frac{\upsilon\rho C_p}{\lambda},
\label{eq:Pr}
\end{equation}
where $\lambda$ is the fluid thermal conductivity. Every case uses the same two-wall averaging and fluid-temperature-reference convention, although the numerical value of $T_{\mathrm{ref}}$ varies with the saved fluid temperatures. This postprocessing definition does not imply thermally periodic temperature profiles.
